\section{Discussion}
\label{sec:discussion}

\subsection{What the audits tell us}
The SPT 2018 and D1 likelihoods do not predict each other well. At the audited fits, parameters chosen by one lens cost the other tens to hundreds of units of $\chi^2$, and the cost depends strongly on direction: testing D1 at a 2018 fit is always the more expensive transfer. Extra D1 multipole coverage, the two releases' different optical-depth priors and default Boltzmann precision change these numbers by a few percent at most. Re-fitting a capped set of nuisance parameters absorbs a large share of the fixed-spectrum disagreement, but not all of it. Where the disagreement sits depends on what else is in the fit. Temperature at $2000\le\ell<3000$ leads both directions with SPT and DESI. Polarization leads the D1-given-2018 direction once the Planck subset and PR4 lensing are added.

Read through the SBT lens, the $\mnu$ shift reported in \cite{arxiv2601_16277} is a change of package under a lens swap, and the two lenses disagree at the level of the packaged likelihoods themselves. Until that disagreement is understood, a shift in a tight limit should be regarded as sensitive to packaging. This conclusion does not depend on the size of the shift, and it does not imply that either release is wrong.

\subsection{Limitations}
\begin{itemize}[leftmargin=*, itemsep=2pt]
  \item \textbf{Optimization.} All audit values are differences between numerically optimized candidates without a global-optimality certificate. Improved source fits moved individual values by up to a few percent. The positivity of $\dchi$ is guaranteed only between the retained candidates.
  \item \textbf{Significance.} The two releases differ in seasons, noise and nuisance models. Without a declared null distribution the $\dchi$ values are comparative diagnostics, not $p$-values. Signed allocations are not independent chi-squares.
  \item \textbf{Nuisance scope.} The SPT+DESI audits and chains fix SPT nuisance parameters at packaged defaults. The full baseline frees only three, and the fixed-spectrum audit re-fits a capped set of twelve. Broader nuisance freedom could change both the audit values and the posteriors.
  \item \textbf{Baseline scope.} Our full baseline omits SPT lensing and differs from the SPT+DESI baseline in Boltzmann code and free foregrounds. Comparisons between the two baselines do not isolate a single cause.
  \item \textbf{Posteriors.} We have no converged $\mnu$ posterior on a qualified numerical target (\cref{sec:results-mnu}). The likelihoods' sensitivity to solver version, caching and precision settings is part of the reason.
  \item \textbf{Released objects only.} The audits operate on released likelihoods. They say where the packaged summaries disagree, not whether calibration, beams, transfer functions, foreground modeling, map-making or covariance construction is responsible.
\end{itemize}

\subsection{Next diagnostics}
The decisive test remains the one proposed in \cite{arxiv2601_16277}: reprocess the 2018 raw data through the D1 pipeline. Short of that, three steps follow from this work.
\begin{enumerate}[leftmargin=*, itemsep=2pt]
  \item \textbf{Predictive audits.} Replace best-fit transfer with posterior predictive checks across lenses \cite{gelman_meng_stern1996_ppc}, with a declared null that accounts for differing noise and seasons. This would turn the localization maps into calibrated statements.
  \item \textbf{Controlled targets.} Fix and record the solver version, precision settings and caching behavior. Verify stored chain rows against fresh evaluations. Only then sample to the quantile precision required by \cref{prop:readout}, given the size of the shift to be resolved.
  \item \textbf{Targeted rewrites.} Where the mismatch localizes (EE at $2000\le\ell<3000$ in the full baseline, TT in the same range with SPT and DESI), test small template or calibration corrections fitted on one lens and evaluated on the other. A correction selected to fit the residual it removes improves the fit by construction. Its value must be judged on held-out data.
\end{enumerate}

\subsection{Practical recommendation}
When a likelihood update moves a tight bound, we suggest reporting, alongside the bound: the directional mismatch in both directions; its signed, covariance-respecting localization by spectrum and multipole, for each baseline used; at least one support control; and posterior diagnostics with quantile-precision targets matched to the size of the claimed shift. These are inexpensive compared with the inference itself, and they change how a moving bound should be read.
